\chapter{Bevezetés}

\section{Problémafelvetés}

A kompetitív online játékok egyik visszatérő problémája, hogy a résztvevők nagyon erős véleményeket fogalmaznak meg egymás teljesítményéről, miközben ehhez rendszerint csak részleges, szubjektív és rövid távú benyomások állnak rendelkezésre. Egy elveszített mérkőzés után gyorsan megjelenik a bűnbakkeresés, de sokkal ritkábban történik meg az, hogy a játékosok valóban adatalapon próbálják megérteni, mi történt a meccsben, hogyan teljesítettek ők maguk, és hogyan helyezhető el ez a teljesítmény egy tágabb közösségi hálózatban.

A \emph{League of Legends} különösen érdekes terep ehhez, mert egyszerre tartalmaz rendkívül gazdag mérkőzésadatokat, ismétlődő csapattársi és ellenfél-kapcsolatokat, valamint olyan társas mintázatokat, amelyek gráfként ábrázolhatók és elemezhetők. Ha egy játékos hosszabb időn keresztül ugyanazokkal az emberekkel kerül egy meccsbe (akár szövetségesként, akár ellenfélként), akkor ebből egy olyan kapcsolati háló építhető fel, amely túlmutat az egyetlen mérkőzés pillanatnyi statisztikáin.

Ebből a felismerésből indult a dolgozat: a cél nem egyszerűen egy újabb meccsnéző vagy statisztikai felület volt, hanem egy olyan rendszer, amely a valós mérkőzéstörténetből játékoshálót épít, ezt a hálót gráfelméleti és rendszerszintű szemlélettel elemzi, majd az eredményeket egy vizuálisan erős, ugyanakkor védhető frontend felületen mutatja meg.

\section{Motiváció}

A projekt mögött részben személyes indíttatás húzódik. Rangsorolt meccsek után visszatérő kérdés volt, hogy a sorozatos negatív találkozások véletlen-e, vagy van mögöttük valami szerkezet -- esetleg ugyanazok az emberek kerülnek újra és újra egymás mellé, csak ezt a matchmaking felülete nem mutatja meg. Ez önmagában persze nem tudományos kérdés, de kiindulópontként elég volt. Ha a Riot API részletes meccsadatokat ad, és ezek ismétlődő játékoskapcsolatokat tartalmaznak, akkor ezekből elvileg gráf építhető -- és ha gráf van, akkor már van mit elemezni.

A rendszer végül jóval nagyobbra nőtt, mint amire számítottam. Ami kezdetben gyors adatvizualizációnak indult, az fokozatosan külön adatgyűjtő réteggé, normalizált pontszámítási pipeline-ná, SQLite-alapú perzisztenciává, Rust futtatókörnyezetté és többoldalas React frontenddé alakult. A Python réteg adatgyűjtési és legacy-kompatibilitási szerepben maradt meg -- ez nem tervezett döntés volt, hanem a fejlesztés természetes következménye.

\section{A dolgozat célja}

A dolgozat célja egy olyan, valós mérkőzésadatokra épülő, többkomponensű rendszer bemutatása, amely:

\begin{itemize}
    \item League of Legends meccsekből játékosok közötti kapcsolatokat gyűjt és tisztít,
    \item a játékosokhoz szerepkör-érzékeny teljesítménymutatókat rendel,
    \item az ismétlődő együtt- és egymás ellen játszásból gráfot épít,
    \item a gráfot klaszterekre, útvonalakra, core-periphery mintázatokra és kutatási célú metrikus kapcsolatokra bontja,
    \item a számításigényes részeket determinisztikus, exportálható és újrafuttatható módon kezeli,
    \item az eredményeket interaktív, de nem UI-első szemlélettel mutatja meg.
\end{itemize}

Fontos szempont volt végig, hogy a rendszer ne csak vizuálisan legyen érdekes, hanem a mögötte lévő döntések is megindokolhatók legyenek: honnan jönnek az adatok, hogyan számolódnak a pontszámok, mit jelent pontosan egy él a gráfban, és mit nem.

\section{Kutatási kérdések}

A jelen dolgozat a következő kutatási kérdésekre keres választ:

\begin{itemize}
    \item Hogyan építhető fel valós League of Legends mérkőzéstörténetből egy olyan játékosháló, amely nem pusztán az együtt játszást, hanem az ellenfél-kapcsolatokat is megőrzi?
    \item Milyen módon lehet a játékosok teljesítményét úgy pontozni, hogy az figyelembe vegye a tényleges szerepkört, de közben megmaradjon értelmezhetőnek?
    \item Mennyiben indokolt a rendszer különböző rétegeit Python, SQLite, Node.js, Rust és React komponensekre bontani?
    \item Mutat-e a jelenlegi gráf pozitív asszortativitást a \texttt{opscore} és \texttt{feedscore} tekintetében?
    \item Miben különbözik a \texttt{flexset} és a \texttt{soloq} kapcsolati szerkezete, és mennyiben írható le a \texttt{flexset} asszociatív core-periphery játékosgráfként?
\end{itemize}

\section{A dolgozat hozzájárulása}

A dolgozat gyakorlati és szakmai hozzájárulásai a jelenlegi implementáció alapján a következők:

\begin{itemize}
    \item több-adatkészletes, regiszter-alapú mérkőzésgyűjtési és feldolgozási architektúra;
    \item szerepkörérzékeny, adatbázis-alapú \texttt{opscore}/\texttt{feedscore} normalizáció;
    \item Rust-központú Graph V2 gráfépítés, klaszterezés és artifact-export SQLite perzisztenciával;
    \item Rust oldali útvonalkeresés BFS, Dijkstra, kétirányú keresés, egzakt A*, asszortativitás és Brandes-központisági támogatással;
    \item \texttt{flexset} és \texttt{soloq} dataset-összehasonlítás asszociatív gráfértelmezéssel;
    \item játékosteljesítményre számolt asszortativitási elemzés több gráfmódban;
    \item előre kiszámított 3D globális gráfnézet és többoldalas, kutatásbarát frontend.
\end{itemize}

\begin{figure}[H]
    \centering
    \includegraphics[width=0.92\textwidth]{match-analysis-page.png}
    \caption{A jelenlegi frontend egyik belépési pontja: meccselemzési és adatgyűjtési felület}
\end{figure}
